创新与伦理

自动驾驶：AI并不与你“共命运”

信息爆炸对认知的影响：成瘾

局部最优解与信息孤岛：大行其道的推荐系统带来思考的惰性、认知的狭窄、……

信息工具无罪，但要慎重圈定使用范围。


如果一个AI工具，同时
拥有存储人类历史上的所有知识，
做数学题
会编程
会艺术创作(诗歌，文章，图片，语音，视频，…)
会行动（配合机器人）
…
人类物种的优势何在？

####比尔盖茨提出三个指导对话的原则：
#####为了充分利用这项非凡技术，应该尝试平衡对人工智能缺点的恐惧与它改善人们生活的能力。既要防范风险，又要让尽可能多的人受益。
#####市场力量不会自然而然地生产出帮助最贫困人群的人工智能产品和服务。相反的可能性更大。需要让世界上最好的人工智能专注于最大的问题。
#####世界需要制定道路规则，让人工智能的任何缺点都远远超过它的好处，让每个人都能享受这些好处，无论他们住在哪里，无论他们有多少钱。

站在人工智能这个巨人的肩膀上，创造更好的人工智能产品和服务，为人类带来生而平等的平和感、富足感。


不管问题多么复杂，如果满足以下五个条件，计算机绝对能够战胜人类。只要有任何一个或者多个条件不满足，计算机做起来就困难了。

1.必须拥有充足的数据或知识；充足不仅指数量大，多样性，不能残缺等。
2.确定性信息；围棋☑️自动驾驶❌
3.完全信息；无论多么复杂，本质上只需要计算速度快。围棋☑️自动驾驶❌
4.静态；按确定性的规律演化，即可预测性。围棋☑️自动驾驶❌
5.特定领域；单任务。围棋☑️自动驾驶（行人、其他车辆）❌








围棋和扑克，哪个容易被AI攻克？博弈游戏



**不确定信息**



**可观测状态信息**



要严密地理解伦理，必须先理解数学。---苏格拉底(公元前469—公元前399)
解决这个问题(在人工智能中编入道德)是一个值得投入下一代最伟大的数学人才去解决的研究挑战。---尼克·博斯特罗姆(1973—)



伊萨克·阿西莫夫在他的科幻小说中提出了一个有关机器人道德的框架，其中列出了多种行动之间的优先顺序。根据这一顺序，机器人应该最优先保护人类不受伤害。只有在满足第一法则的前提下，机器人才应该遵循人类的指令。然而，指望这样一张优先列表能够解决所有与道德决策相关的问题简直是天方夜谭。这是因为应该采取的行动不仅取决于所有可能的行动，还取决于在什么情景下进行决策。但是可以想象的情景数量众多，并且呈指数增长。列出在所有情况下可以采取的所有行动，就相当于写出一个算法，能够在所有情境下判断出应该执行什么行动。但我们基本上可以确定，这样的算法会拥有巨大的所罗门诺夫复杂度。无论是对人类还是对机器而言，这种算法的编写、读取和应用完全不切实际。

知识是合理的目的吗?
现在剩下的工作就是确定什么结果才算可取。我们对社会真正的期望是什么?文明的最终目标是什么?我们在客观上有什么作用?这就是结果论者要回答的根本问题。
科学工作者经常提倡对知识的渴望。将知识作为社会的目标，这在道德上成立吗?要求大多数人拥有知识，这是否的确合理?应不应该要求所有人都吞下那颗红色药丸?大概不行。实际上，即使只是将知识定为科学的目标，也会导致许多奇怪的结果。
问题在于世界既庞大又复杂。这个世界中可以知晓的事物的数量远远超越了我们的认知能力。欧内斯特·卢瑟福甚至主张“所有科学要么是物理，要么就是集邮”。庞加莱也这样说：“事实的堆砌与一门科学的距离并不比一堆石头与一座房屋的距离近。”我个人也没有任何耐心去背诵那些由零碎知识组成的列表。
与其只关心数据，人们更希望寻求这些数据中暗含的理论，常用的方法就是应用贝叶斯公式。但这样的话，目标是什么?就是进行预测吗?我们可不可以认为尝试得出正确结论就是合理的目的?答案绝对是否定的。即使是在KL散度这种精细的意义上，尝试得出正确结论也有局限性。


